% LSPV-4 REPORT::TRACEISO1::2026-10-14T04:02:00+09:00
% !TEX program = pdflatex
\documentclass[11pt]{article}
\usepackage[T1]{fontenc}
\usepackage{lmodern,amsmath,amssymb,amsthm,booktabs,tabularx,xurl,hyperref}
\usepackage[margin=0.80in]{geometry}
\hypersetup{hidelinks}
\setlength{\emergencystretch}{3em}
\newtheorem{theorem}{Theorem}
\newtheorem{corollary}[theorem]{Corollary}
\newcommand{\Tr}{\operatorname{Tr}}
\newcommand{\F}{\mathbb F}
\title{TRACEISO1: Deterministic Trace Root Isolation\\\large All Explicit Binary-Field Cube Residuals and a Constructive Isomorphism Bridge}
\author{UNIVERSAL\_CORE\_BRIDGE / LSPV-4}
\date{14 October 2026, 04:02 KST}
\begin{document}\maketitle
\paragraph{Status.}
\textbf{CANDIDATE integrated bridge; complete proof, independent specialist review pending.}
Finite-field factorization and isomorphism are established algorithmic subjects,
not new foundational discoveries of this run. The new project contribution is
an explicit, tested characteristic-two selector and its charged Tseitin interface.
It closes the previously stated arbitrary-polynomial-basis cube-selection gap
for every degree, not only the cyclotomic sequence. It does not recover hidden
field structure from arbitrary quadratic equations. \textbf{P versus NP remains OPEN.}

\section{Target, representation, and literature correction}
The starting question was whether the CYCLOROOTFW1 selector on
$d=2\cdot3^a$ can be transported to arbitrary irreducible presentations.
Lenstra's deterministic polynomial-time finite-field isomorphism theorem is
already established \cite{lenstra}; its statement and explicit-data convention
are confirmed in the later primary research paper of Brieulle et al.
\cite{brieulle}. Berlekamp's factorization work provides the classical
linear-algebra/GCD background \cite{berlekamp}. We give a self-contained
specialized proof rather than assuming a black-box isomorphism algorithm.
The original Lenstra proof was not retrieved in full; this report does not
claim to have re-audited that proof.

A stronger route bypasses the isomorphism entirely. The field is
$K=\F_2[X]/(f)$, where the input explicitly lists all $d+1$ coefficients of
an irreducible degree-$d$ polynomial. Elements have $d$ coefficient bits.
Our input is not a succinct exponent description, unknown multiplication
oracle, or anonymized quadratic ANF. Irreducibility and types are checked.
No source root, cubic nonresidue, randomness, or nonuniform advice is supplied.

\section{One concrete lemma: deterministic trace separation}
\begin{theorem}[Explicit binary-field root isolation]\label{thm:root}
One deterministic uniform algorithm, given $d\ge1$, irreducible
$f\in\F_2[X]$ of degree $d$, and a densely represented univariate
$p(T)\in K[T]$ of degree at most $m$, returns a root in $K$ if one exists
and otherwise reports no root. Its time and workspace are polynomial in
$d+m$ and the explicit coefficient length.
\end{theorem}
\begin{proof}
Handle the zero polynomial by returning zero, and a nonzero constant by
rejecting. Otherwise replace $p$ by its monic associate. Compute
\[
 g=\gcd\bigl(p,T^{2^d}-T\bigr)                         \tag{1}
\]
using $d$ successive squarings modulo $p$, not an expanded polynomial of
exponential degree. Since $T^{2^d}-T$ has every element of $K$ as a simple
root, $g$ is exactly the squarefree product of all field-linear factors of
$p$. If $g=1$, there is no field root.

Use the public basis $\beta_j=X^j\bmod f$, $0\le j<d$. For each $j$ while
$\deg g>1$, compute
\[
 h_j(T)=\sum_{i=0}^{d-1}(\beta_j T)^{2^i}\bmod g,
 \qquad g_0=\gcd(g,h_j).                              \tag{2}
\]
Again only modular squarings are used. At a root $r$ of $g$, $h_j(r)$ equals
$\Tr_{K/\F_2}(\beta_jr)\in\{0,1\}$. If $g_0$ is nonconstant retain
$g_0$; otherwise retain $g$, whose remaining roots all have trace coordinate
one. The retained root set is never empty, and all retained roots agree on
every coordinate processed so far.

The pairing $(s,t)\mapsto\Tr(st)$ is nondegenerate. Indeed, the trace
polynomial $Z+Z^2+\cdots+Z^{2^{d-1}}$ is nonzero and has degree less than
$2^d$, so it cannot vanish on all of $K$. Multiplication by any nonzero
$\delta$ permutes $K$. Therefore $\Tr(\beta_j\delta)$ cannot vanish for
all members of a basis. Taking $\delta=r+s$ shows that any two distinct
roots differ in at least one of the $d$ trace coordinates.

After processing the basis, the retained nonempty set consequently contains
one root. The final factor is $T+r$. Return $r$ and verify $p(r)=0$ against
the original coefficients. This proves existence, correctness, termination,
and rejection of rootless inputs without enumerating $K$.
\end{proof}

\paragraph{Charged cost.}
Naive polynomial arithmetic uses $O(d^2(m+1)^2)$ field operations for
(1)--(2), including polynomial GCD work. Polynomial-basis field arithmetic
and inversion use polynomial bit time; irreducibility is checked by
successive Frobenius powers and binary polynomial GCDs. A deliberately loose
multitape bound allowing linear scans for coefficient access is
\[
 O\bigl(d^5(m+1)^3+d^4\bigr)\text{ bit time},
 \qquad O\bigl((d+m+1)^3\bigr)\text{ writable bits}.     \tag{3}
\]
These are conservative algorithmic upper bounds, not measured instruction
counts or optimality claims. Counters and temporary polynomials are included.
The audit's 25-second deadline applies only to test invocations; imported
production routines have no fixed timeout and no fixed degree cap.

\begin{corollary}[All-degree target-only cube selector]\label{cor:cube}
For every $d\ge1$, every explicit irreducible degree-$d$ binary modulus,
and every $b\in K$, a single deterministic polynomial-time program returns
$z$ satisfying $z^3=b$ exactly when one exists.
\end{corollary}
\begin{proof}
Apply Theorem~\ref{thm:root} to $p(T)=T^3+b$. Equation (1) also handles
$b=0$ and repeated roots correctly. The degree $m=3$ is constant in (3).
\end{proof}

Thus fixed valuation, a logarithmic subfield window, a public nonresidue,
and the canonical cyclotomic modulus are not necessary for polynomial
cube selection in this explicit characteristic-two model. This supersedes
only the corresponding open algorithmic regime, not any earlier valid
selector or restricted proof-system theorem.

\section{The originally requested representation bridge}
\begin{corollary}[Constructive binary-field isomorphism]\label{cor:iso}
Given two irreducible degree-$d$ binary moduli $f,g$, one can compute
forward and inverse $\F_2$-linear field-isomorphism matrices in deterministic
polynomial time.
\end{corollary}
\begin{proof}
In $K_g=\F_2[Y]/(g)$, the polynomial $f(T)$ splits completely: its degree-$d$
roots lie in the unique field with $2^d$ elements. Find one root $\alpha$
by Theorem~\ref{thm:root}. Evaluation gives
\[
 \varphi\left(\sum_{i=0}^{d-1}c_iX^i\right)
       =\sum_{i=0}^{d-1}c_i\alpha^i.                    \tag{4}
\]
Irreducibility makes $1,\alpha,\ldots,\alpha^{d-1}$ a basis, so (4) is a
field isomorphism. These powers are the columns of a binary $d\times d$
matrix. Gaussian elimination constructs its inverse. Matrix generation,
storage, evaluation and inverse evaluation are all charged. Substituting
$m=d$ in (3) gives the coarse $O(d^8)$ bit-time bound; matrices use $O(d^2)$ bits.
\end{proof}

For $d=2\cdot3^a$, take $g(Y)=Y^d+Y^{d/2}+1$ from CYCLOROOTFW1.
The implemented comparison transports a target into this canonical field,
uses its separately implemented correction selector, and transports the
root back. Direct and transported selectors need not return the same root;
both must return a valid root or reject a noncube. The direct trace selector
needs no cyclotomic representation at all.

\section{Coupled Tseitin family and witness maps}
Let $X=(V,E)$ be an explicitly listed simple graph, retaining isolated
vertices. Fix a deterministic spanning forest $T$ and ordered cotree
$C=E\setminus T$. Put $d=|C|\ge1$. For a charge $c\in\F_2^V$ and a field
target $b\in K$, define
\[
 G_{X,f,c,b}(x)=[B_Xx=c]\ \land\ [(\pi_Cx)^3=b].        \tag{5}
\]
The record checks vertex/edge types, endpoints, uniqueness, degree agreement,
field irreducibility and component charge parity. Every component, including
an isolate, must have even total charge. Malformed records are distinct from
well-formed NO instances with inconsistent charge or a noncube target.
The executable raises a validation error on malformed records; a total
language decider maps that error to rejection.

Set $a=\operatorname{Lift}_T(c,0)$ by leaf-to-root forest solution. For each
cotree coordinate $j$, set $L_j=\operatorname{Lift}_T(0,e_j)$. Then
\[
 B_Xa=c,\qquad B_XL=0,\qquad \pi_CL=I_d,
 \qquad x=a+Lz.                                       \tag{6}
\]
The residual in (5) is precisely $z^3=b$. Select $z$ using
Corollary~\ref{cor:cube}, then reconstruct $x$ with the same forest. Projection
and lifting are inverse witness maps between residual roots and full models
of (5). Enumeration of the full solution set is unnecessary. Materializing
$L$ would use $O(|E|d)$ bits; the implementation instead lifts directly and
also explicitly audits basis columns and $a+Lz$ on small instances.

The explicit record length is polynomial in
$(|V|+|E|)\log(|V|+|E|+2)+d$. Forest discovery, component testing, selection,
reconstruction, original-equation verification and output are all polynomial
in that length. Boolean cubing is quadratic because squaring is linear over
$\F_2$ and multiplication is bilinear. Its auxiliary Boolean circuit has
polynomial size; encoding it does not license recognition of arbitrary
untagged SAT presentations as this typed family.

For a general residual support $R$, retain the existing exact identity
\[
 \rho=\dim\pi_R\ker B_X
       =|R|+\operatorname{cc}(X)-\operatorname{cc}(X-R). \tag{7}
\]
The kernel of this projection is the cycle space of $X-R$, which proves
(7) by subtracting cycle-space dimensions, with all isolates included.
For $R=C$, $\rho=d$. The familiar $2^\rho$ enumeration is an upper bound;
(5) has a polynomial selector despite full projected rank.

\paragraph{Source compatibility limitation.}
With a source conjunct $F(u)$ and circuits $b(u),c(u)$, fixing a supplied
$u$ lets us evaluate the circuits, test parity and cube-image membership,
and select an extension in polynomial time. This does not find $u\models F$
and does not guarantee that every source witness is compatible: noncube
$b(u)$ can destroy satisfiability. For odd $d$ all targets are cubes; for even
$d$ the nonzero cube image is proper. No generic amplification follows.

\section{Actual bounded experiments}
Nineteen completed, separately checkpointed invocations made
\textbf{26,359 assertions with zero failures}. All outputs pin the code
SHA-256 and test seed. The slow independent carryless multiplication routine
checks the exhaustive cube truth tables; a second checker verifies coupled
witnesses using both original parity equations and $a+Lz$. These are separate
checks by the same author, not independent specialist review.
\begin{center}\small
\begin{tabularx}{\linewidth}{p{0.30\linewidth}Xr}\toprule
Batch group & Coverage & Assertions\\\midrule
Exhaustive binary fields & All 71 monic irreducible presentations at degrees
1--8; all 10,824 targets & 16,138\\
Explicit isomorphisms & Seven presentations at degrees 6, 18, 54; forward/inverse
basis tests and every basis-product pair & 7,696\\
Larger direct selectors & Degrees 12, 31, 36, 100, 162, 486; seeded targets,
noncubes, forest witnesses & 1,052\\
Structural/typing controls & 729 projected-rank cases with isolates, 40 public
affine conjugations, source and polynomial controls & 952\\
Separate coupled checker & 192 target/charge cases, 66 positive witnesses,
128 inconsistent-charge rejections & 521\\\bottomrule
\end{tabularx}\end{center}
The largest completed batch took 4.645 seconds. The largest isomorphism
actually tested has degree 54; the degree-486 tests are direct cube selection,
not an isomorphism benchmark. There were 376 positive complete witness
reconstructions across the larger/isomorphism batches and the separate checker.
Ten malformed field records, four malformed graph cases and three malformed
coupled records were rejected. Source controls include Horn, Krom, affine,
independent-product, general 3-CNF and UNSAT cases. Production selection is
deterministic; random generation of large test moduli and targets is not part
of its claimed algorithm. No exhaustive large-field scan is performed.

One shell command grouping several early invocations exceeded its outer tool
timeout. Completed outputs were retained and the remaining batches rerun
separately. No final completed invocation exceeded its cooperative 25-second
budget. The draft duplicate test is excluded from the final totals and package.
Operation counters are diagnostic, not full bit-operation or memory profiles.
Finite results are \textbf{COMPUTATIONAL}; the asymptotic argument is the proof above.

\section{Corpus, integration, and barrier audit}
\paragraph{CORPUS REVIEWED.}
Read current AGENTS.md, CODEX.md, the repository research skill, the redirect,
root main.tex, the current recursive tree, actual main4/main5 driver portions,
CYCLOROOTFW1, VALSUBCUBEFW1 and MAIN VALSUBREVIEW1. Read the predecessor's
required header/body/final sentinel; its latest located directory artifact
is VALSUBCUBEFW1. Searched the corpus for the cyclotomic/valuation claims and
trace splitting. Earlier task and EARLY22/SINGLE22 sources remain in the
conversation record; this run did not newly reread every task, historic
manuscript include, original Drive passage or exploration file. No complete
local checkout was available because local GitHub DNS resolution failed.
Repository-wide preflight/validator execution is not claimed.

\paragraph{EARLY22/SINGLE22 acknowledgment.}
\textbf{ALREADY INTEGRATED}, via their retained ledger in root main.tex and
IMAGE22COMP1. Existing signed exact-$(2,2)$ compilers, structured minimality,
Resolution/PC projections and the separate clause-monotone construction keep
their hypotheses; historical variable counts are not encoded bit lengths.
The affine continuation limitation remains attached. No original missing
scientific body is reconstructed. Exact new location: Sections 2--5 here and
the proposed LSPV-4 root block. Next check is raw-ANF recognition, not another
punctured-affine or proof-size-to-runtime promotion.

\paragraph{DUPLICATION / CONTRADICTION AUDIT.}
General finite-field algorithms are known literature, not novel project
claims. No matching all-degree trace selector was found in the bounded
project search. Earlier nonresidue/subfield selectors remain valid specialized
algorithms. Their statement that growing valuation or arbitrary explicit
binary modulus leaves cube selection open is superseded by Theorem~\ref{thm:root}.
An opaque multiplication structure is a different input model.

\paragraph{WHAT CHANGED RELATIVE TO ROOT MAIN.}
The bridge no longer needs fixed valuation, a subfield enumeration window,
a canonical modulus or a supplied nonresidue. It has a direct all-degree
selector and constructive forward/inverse representation maps. Only the
LSPV-4 ownership region is authorized for reconciliation. Stale executive or
other-specialist regime descriptions are a synthesis handoff, not permission
to edit their bytes.

\paragraph{ADOPTED / REJECTED / RETIRED PRIOR COMPONENTS.}
Adopted: affine elimination, projected-rank identity, strict types, forest
lifting, public field arithmetic and existing exact proof projections.
Rejected as unproved: arbitrary hidden-field recovery and universal source
compatibility. The proposed explicit-field cube family as an unavoidable
superpolynomial-work source is \textbf{FALSIFIED/RETIRED}, subject to review
of this integrated proof. P versus NP and the general nonlinear residual
lower bound remain \textbf{OPEN}.

\paragraph{Models and transfers.}
This is an exact transfer of finite-field root algorithms to their actual
residual object. It is not a transfer from DNNF size or regular Resolution
\cite{colnet} to witness-search time. Arbitrary polynomial workspace is
allowed, separately from constant program length. With the declared decoder,
the uniform selector gives an eventual eligible short-program polynomial
upper bound on this typed UWI subfamily, not on all satisfiable formulas.
Public affine conjugations can be inverted and charged; hidden conjugations
cannot be assumed available. Added semantic or algebraic oracles only help
this algorithm; it supplies no nonrelativizing separation step. No broad
constructive circuit property, ETH, GRH, cryptographic or random-oracle
hypothesis is used. The coding/finite-field transfer is exact; generic PQ
binding, collision entropy, Schur-square or syzygy language supplies no
additional lower bound here.

\section{One next target}
Determine whether a field-cube representation can be recovered from the
explicit quadratic ANF when neither the field modulus nor the invertible
coordinate maps are supplied. State the recognizable domain and charge any
recovery before invoking the selector. A failure of that particular recovery
algorithm would not itself establish a runtime lower bound.

\begin{thebibliography}{9}
\bibitem{lenstra} H. W. Lenstra, Jr. Finding isomorphisms between finite
fields. Mathematics of Computation 56(193), 329--347, 1991.
DOI: 10.1090/S0025-5718-1991-1052099-2.
\bibitem{brieulle} L. Brieulle, L. De Feo, J. Doliskani, J.-P. Flori,
E. Schost. Computing isomorphisms and embeddings of finite fields.
\url{https://arxiv.org/abs/1705.01221}.
\bibitem{berlekamp} E. R. Berlekamp. Factoring polynomials over finite
fields. Bell System Technical Journal 46(8), 1853--1859, 1967.
DOI: 10.1002/j.1538-7305.1967.tb03174.x.
\bibitem{colnet} A. de Colnet and S. Mengel. Characterizing Tseitin-formulas
with short regular resolution refutations. JAIR 76, 265--286, 2023.
\url{https://arxiv.org/abs/2103.09609}.
\end{thebibliography}
\end{document}
% END_OF_VERIFIED_REPORT::LSPV-4::TRACEISO1
